
\documentclass{article}
\usepackage{colortbl}
\usepackage{makecell}
\usepackage{multirow}
\usepackage{supertabular}

\begin{document}

\newcounter{utterance}

\centering \large Interaction Transcript for game `cladder', experiment `full\_v1.5\_default', episode 5155 with reflexion{-}llama{-}3.1{-}8b{-}instruct.
\vspace{24pt}

{ \footnotesize  \setcounter{utterance}{1}
\setlength{\tabcolsep}{0pt}
\begin{supertabular}{c@{$\;$}|p{.15\linewidth}@{}p{.15\linewidth}p{.15\linewidth}p{.15\linewidth}p{.15\linewidth}p{.15\linewidth}}
    \# & \multicolumn{2}{c}{Player} && \multicolumn{2}{c}{Game Master} \\
    \hline

    \theutterance \stepcounter{utterance}  
    & & & \multicolumn{4}{p{0.6\linewidth}}{
        \cellcolor[rgb]{0.9,0.9,0.9}{
            \makecell[{{p{\linewidth}}}]{
                \texttt{\tiny{[P1$\langle$GM]}}
                \texttt{You are an expert in causal inference. The following question is not a typical commonsense query, but rather a meticulously designed question created by a professor specializing in causal inference, intended to assess the students' mastery of the course content.} \\
\texttt{Imagine a self{-}contained, hypothetical world with only the following conditions, and without any unmentioned factors or causal relationships: Talent has a direct effect on elite institution admission status. Brown eyes has a direct effect on elite institution admission status. Method 1: We look at how talent correlates with brown eyes case by case according to elite institution admission status. Method 2: We look directly at how talent correlates with brown eyes in general. To understand how talent affects brown eyes, is it more correct to use the Method 1 than Method 2?} \\
\texttt{Start your answer with 'Yes' or 'No', followed by additional reasoning or evidence to support your explanation.} \\
            }
        }
    }
     \\ \\

    \theutterance \stepcounter{utterance}  
    & \multicolumn{4}{p{0.6\linewidth}}{
        \cellcolor[rgb]{0.561,0.737,0.561}{
            \makecell[{{p{\linewidth}}}]{
                \texttt{\tiny{[P1$\rangle$GM]}}
                \texttt{No, it is not more correct to use Method 1 than Method 2 to understand how talent affects brown eyes.} \\
\\ 
\texttt{In this hypothetical world, talent and brown eyes both have direct effects on elite institution admission status. This creates a confounding variable problem, where the relationship between talent and brown eyes is confounded by their shared effect on elite institution admission status.} \\
\\ 
\texttt{Method 1, which examines the correlation between talent and brown eyes case by case according to elite institution admission status, is essentially a form of stratification or matching. However, this method does not address the underlying causal relationship between talent and brown eyes. Instead, it only examines the correlation between the two variables within each stratum of elite institution admission status.} \\
\\ 
\texttt{Method 2, which directly examines the correlation between talent and brown eyes in general, is a more straightforward approach. However, in the presence of a confounding variable (elite institution admission status), this method may also be subject to bias.} \\
\\ 
\texttt{To accurately understand the causal relationship between talent and brown eyes, we would need to use a method that controls for the confounding variable, such as instrumental variables (IV) or regression adjustment. However, given the hypothetical scenario, we can use the concept of causal diagrams or directed acyclic graphs (DAGs) to visualize the relationships.} \\
\\ 
\texttt{In this case, the DAG would show that talent and brown eyes both have direct effects on elite institution admission status, and that elite institution admission status is a common cause of both talent and brown eyes. This would indicate that the relationship between talent and brown eyes is confounded by elite institution admission status.} \\
\\ 
\texttt{Therefore, neither Method 1 nor Method 2 is sufficient to accurately understand the causal relationship between talent and brown eyes. A more sophisticated method, such as IV or regression adjustment, would be needed to control for the confounding variable and estimate the causal effect of talent on brown eyes.} \\
            }
        }
    }
    & & \\ \\

    \theutterance \stepcounter{utterance}  
    & & & \multicolumn{2}{p{0.3\linewidth}}{
        \cellcolor[rgb]{0.9,0.9,0.9}{
            \makecell[{{p{\linewidth}}}]{
                \texttt{\tiny{[GM$|$GM]}}
                \texttt{No} \\
            }
        }
    }
    & & \\ \\

    \theutterance \stepcounter{utterance}  
    & & & \multicolumn{2}{p{0.3\linewidth}}{
        \cellcolor[rgb]{0.9,0.9,0.9}{
            \makecell[{{p{\linewidth}}}]{
                \texttt{\tiny{[GM$|$GM]}}
                \texttt{no} \\
            }
        }
    }
    & & \\ \\

    \theutterance \stepcounter{utterance}  
    & & & \multicolumn{2}{p{0.3\linewidth}}{
        \cellcolor[rgb]{0.9,0.9,0.9}{
            \makecell[{{p{\linewidth}}}]{
                \texttt{\tiny{[GM$|$GM]}}
                \texttt{game\_result = WIN} \\
            }
        }
    }
    & & \\ \\

\end{supertabular}
}

\end{document}
